\section{Discussion, Limitations, and Extensions}
\label{sec:discussion}

\subsection{Positioning of the formalism}
\label{subsec:disc-positioning}

\Cref{sec:formal-model,sec:session-persistence,sec:theoretical} together specify the object, the pipeline, and the epistemic status of deferred claims.
At a high level, State of Thought answers a structural question---what counts as \emph{the} cognitive state of a multi-agent stack---without tying the answer to a vendor API or a particular retrieval implementation.
The primitive commitments are threefold: a \emph{global} subgraph of closed, durable thoughts; \emph{session} slices that make multi-agent expansion finite and inspectable; and an \emph{external} interface that lands evidence in layer~$0$ before abstraction.

\subsection{Layers, slices, and closure}
\label{subsec:disc-mechanisms}

\paragraph{Layers.}
Layer indices realize the vertex partition $\Vlayer{h}$ and enforce admissibility: higher-layer thoughts that close globally must declare downward witness structure to globally persistent layer~$0$ (\Cref{def:admissibility}).
Summarization is therefore paid for explicitly: it cannot silently float free of evidence.

\paragraph{Slices.}
Retrieval selects a finite working substructure from a potentially unbounded global graph; without slices, parallel expansion is not a well-defined finite object.

\paragraph{Closure.}
Persistence is gated exclusively by $\CloseOp_t^{\Pi}$ (\Cref{def:closure-op}): admissibility gate, conflict policy, then global write.
Merge reconciles agents; it does not, by itself, assert durability.

\subsection{Reference policies as illustration}
\label{subsec:disc-reference}

The triple $(\PsiConf^{\mathrm{ref}},\PsiMerge^{\mathrm{ref}},\Pi_{\mathrm{prov}}^{\mathrm{ref}})$ is a \emph{pedagogical} choice: it exhibits conservative coexistence under symmetric conflict, commutative merge, and trace-minimal provenance without declaring uniqueness or optimality.
Alternative $\PsiConf$ and $\PsiMerge$ remain compatible with the same signatures (\Cref{def:conflict-policy,def:merge-policy}).

\subsection{Limitations and scope of claims}
\label{sec:limitations}

This work is a pure theory-and-method contribution: definitions, operators, and explicitly qualified statements (remarks, one conjecture, and open questions listed in the appendix).
It does not describe an implementation, measure a runtime, or compare systems on tasks; by design it supports no quantitative claim about deployment behavior, scalability, latency, or task success.

\paragraph{Policy underdetermination.}
The reference triple fixes one coherent realization; $\SliceOp_t$, alternative merge orders, and alternative conflict regimes are not pinned to a unique global specification---many realizations fit \Cref{sec:formal-model,sec:session-persistence}.

\paragraph{Continuations of the theory program.}
Admissibility is defined path-wise (\Cref{def:admissibility}).
Refining or resolving \Cref{conj:close-preserve}, and characterizing totality or uniqueness of $\MergeOp_t^{\PsiMerge}$ and $\CloseOp_t^{\Pi}$ for general $\Pi$, are natural mathematical continuations; \Cref{sec:th-open} and \Cref{app:open-questions} collect them in one place.
A lemma--theorem layer on witness-path dynamics under closure would spell out graph-theoretic hypotheses (edge updates under $\CloseOp_t^{\Pi}$, preservation of $\Tadm$-connectivity) beyond this manuscript's scope; see \Cref{app:proof-program}.

\paragraph{Extensions.}
Higher-order topological refinements of the baseline graph model are left for future work.
Comparative empirical studies and alternative policy instantiations beyond the reference triple are outside the scope of this paper; \Cref{sec:related} and \Cref{app:notes} indicate directions.
